%% kapa: exam=1 round=23 year=2012 period=2 form=A subject=appraisal_law subject_ko=부동산관계법규 seq=12 number=52
%% source: Q-Net 감정평가사 기출문제 (cst00309) artlSeq=5052580 file=raw/1008/5052580/2012년 제23회 감정평가사 제1차시험 기출문제(2교시 A형).pdf
%% conversion: kordoc 4.24.51 -> tools/gen1.py (auto); review=proofread-sonnet; flags=none
\begin{question}{52}
국토의 계획 및 이용에 관한 법령상 기반시설설치비용에 관한 설명으로 옳은 것은?
\begin{choices}
  \item 특별시장ㆍ광역시장ㆍ시장 또는 군수는 기반시설설치비용을 부과하려면 부과기준시점부터 30일 이내에 납부의무자에게 적용되는 부과기준 및 부과될 기반시설설치비용을 미리 알려야 한다.
  \item 기존 건축물을 철거하고 신축하는 경우 기존 건축물의 연면적은 고려하지 않고 신축하는 건축물의 연면적만을 고려하여 기반시설설치비용을 부과한다.
  \item 납부의무자는 건축허가를 받은 날부터 2개월 이내에 기반시설설치비용을 납부하여야 한다.
  \item 기반시설설치비용 중 기반시설을 설치하는 데 필요한 기반시설 표준시설비용에는 해당 구역의 평균 개별공시지가가 반영된다.
  \item 납부한 기반시설설치비용은 해당 기반시설부담구역에 필요한 기반시설 및 용지를 모두 확보한 후 잔액이 생기면 납부의무자에게 환급된다.
\end{choices}
\end{question}
\answer{1}
